\chapter*{Περίληψη}
\addcontentsline{toc}{chapter}{Περίληψη}

Η ενεργειακή αναβάθμιση του κτιριακού αποθέματος αποτελεί βασικό άξονα της
ευρωπαϊκής και εθνικής ενεργειακής πολιτικής, καθώς τα κτίρια συνδέονται με
σημαντικό μέρος της τελικής κατανάλωσης ενέργειας και των εκπομπών αερίων του
θερμοκηπίου. Παρά την ύπαρξη δημόσιων προγραμμάτων, τραπεζικών προϊόντων και
ιδιωτικών μηχανισμών χρηματοδότησης, η πληροφορία που αφορά τους όρους, τα
κριτήρια επιλεξιμότητας και τα βασικά χαρακτηριστικά των εργαλείων αυτών
παραμένει συχνά κατακερματισμένη σε διαφορετικές ιστοσελίδες, οδηγούς,
ανακοινώσεις και αρχεία. Η κατάσταση αυτή δυσχεραίνει τόσο τον τελικό χρήστη
όσο και τον μηχανικό ή διαχειριστή που χρειάζεται να συγκρίνει και να
επικαιροποιήσει δεδομένα.

Αντικείμενο της παρούσας διπλωματικής εργασίας είναι η ανάπτυξη μεθοδολογίας
και εφαρμογής για την αυτόματη εξαγωγή και τυποποίηση δεδομένων που αφορούν
χρηματοδοτικά εργαλεία ενεργειακής αναβάθμισης στην Ελλάδα. Η προτεινόμενη
προσέγγιση συνδυάζει τεχνικές διαδικτυακής συλλογής δεδομένων με μεγάλα
γλωσσικά μοντέλα, ώστε να υποστηρίζεται η απόφαση συνάφειας μιας πηγής, η
εξαγωγή δομημένης πληροφορίας, η αποθήκευση σε κοινό σχήμα και ο έλεγχος
ποιότητας των αποτελεσμάτων. Ιδιαίτερη έμφαση δίνεται στη σύνδεση κάθε
εξαγόμενου στοιχείου με την πρωτογενή πηγή, στην αποφυγή μη τεκμηριωμένων
συμπερασμάτων και στην ανάγκη ανθρώπινης επιβεβαίωσης σε πεδία με οικονομικές
και διοικητικές συνέπειες.

Η εργασία στοχεύει να συμβάλει στη μείωση του κατακερματισμού της πληροφορίας
και στην καλύτερη αξιοποίηση των διαθέσιμων χρηματοδοτικών εργαλείων,
αναδεικνύοντας παράλληλα τις δυνατότητες και τους περιορισμούς των μεγάλων
γλωσσικών μοντέλων σε εφαρμογές ενεργειακής πολιτικής όπου η ακρίβεια, η
ιχνηλασιμότητα και η επικαιροποίηση είναι κρίσιμες απαιτήσεις.

\vspace{0.8cm}
\noindent\textbf{Λέξεις-κλειδιά:} ενεργειακή αναβάθμιση κτιρίων, χρηματοδοτικά
εργαλεία, μεγάλα γλωσσικά μοντέλα, διαδικτυακή συλλογή δεδομένων, εξαγωγή
πληροφορίας, δομημένα δεδομένα, ενεργειακή πολιτική.
\clearpage
